\chapter{Построение интеллектуальной системы выбора скорости движения}


\section{Общая схема решения}

Итоговая система состоит из двух уровней.
Нижний уровень представляет собой аналитический регулятор траекторного слежения, который получает текущее состояние квадрокоптера, находит ближайшую точку опорной траектории и формирует управляющие воздействия.
Верхний уровень получает компактный набор признаков, характеризующих локальную ситуацию, и предсказывает целевую параметрическую скорость $V^{\ast}$, с которой объект должен продвигаться вдоль траектории.

Такое разделение ролей выбрано намеренно.
Аналитический контур отвечает за устойчивость и геометрическую корректность движения, а обучаемая модель решает более узкую задачу выбора темпа движения.
За счет этого удается уменьшить риск расходимости системы: даже при неидеальном прогнозе скорости базовый регулятор продолжает удерживать объект около траектории.


\section{Математическая модель квадрокоптера}

В программной реализации используется нелинейная модель квадрокоптера, восходящая к схеме, описанной в диссертационной работе С.~А.~Кима~\cite{kim_algorithms}.
Состояние системы включает положение, линейные скорости, углы ориентации, угловые скорости и дополнительные интеграторные переменные.
Нелинейная динамика квадрокоптера описывается в связанной и инерциальной системах координат.
Вектор положения центра масс обозначим как $p = [x, y, z]^T$, вектор линейной скорости~--- как $v = \dot p$, углы ориентации~--- как $\eta = [\phi, \theta, \psi]^T$, а угловую скорость
в связанной системе координат~--- как $\omega = [p_\omega, q_\omega, r_\omega]^T$.
Тогда поступательное движение квадрокоптера можно записать в виде:
\begin{gather}
    m\dot v = mg e_3 - T R(\eta)e_3 + F_d,
\end{gather}
где $m$~--- масса аппарата, $g$~--- ускорение свободного падения,
$e_3 = [0,0,1]^T$, $T$~--- суммарная тяга винтов, $R(\eta)$~---
матрица поворота из связанной системы координат в инерциальную, а $F_d$~--- вектор внешних возмущений и аэродинамического сопротивления.

Вращательное движение задаётся уравнением Эйлера:
\begin{gather}
    J\dot \omega = \tau - \omega \times J\omega,
\end{gather}
где $J$~--- матрица моментов инерции квадрокоптера,
$\tau = [\tau_\phi,\tau_\theta,\tau_\psi]^T$~--- вектор управляющих моментов
по крену, тангажу и рысканью. Связь между угловой скоростью $\omega$
и производными углов Эйлера задаётся кинематическим соотношением:
\begin{gather}
    \dot \eta = W(\eta)\omega,
\end{gather}
где $W(\eta)$~--- матрица преобразования угловых скоростей.

Таким образом, полное состояние объекта можно представить как
\begin{gather}
    x = [p^T, v^T, \eta^T, \omega^T]^T.
\end{gather}
Данная модель является нелинейной, поскольку матрица поворота $R(\eta)$,
кинематическая матрица $W(\eta)$, а также гироскопический член
$\omega \times J\omega$ зависят от текущего состояния объекта.
Именно поэтому задача траекторного управления не сводится к линейному слежению за координатами и требует отдельного учёта геометрии траектории и динамики движения.
Практический смысл такого расширения в том, что модель остается физически интерпретируемой, но при этом лучше подходит для синтеза и численного моделирования регулятора.

Нейросеть в этой работе не взаимодействует с большим состоянием модели.
Полная динамика квадрокоптера слишком избыточна для задачи выбора скорости.
Поэтому между моделью объекта и ML-модулем вводится слой инженерных признаков, где сложное состояние сжимается до небольшого вектора понятных характеристик.


\section{Геометрия траектории и вычисление ошибок}

Заданная траектория рассматривается как параметрическая пространственная кривая.
Для каждой точки движения необходимо определить ближайшую точку на этой кривой, а затем выразить отклонение квадрокоптера в локальной системе координат, связанной с траекторией.
В работе для достижения поставленных задач рассматриваются прямые, окружности и спирали, а также реализованы итерационные процедуры поиска ближайшей точки.

После этого вычисляются основные ошибки слежения.
Продольная ошибка показывает рассинхронизацию по продвижению вдоль маршрута.
Поперечная ошибка показывает, насколько объект ушел в сторону от траектории.
Дополнительно учитывается ошибка по курсу или рассогласование направления движения с касательной к кривой.
Именно эти величины непосредственно связаны с качеством движения и поэтому используются как при работе регулятора, так и при построении признаков для моделей машинного обучения.


\section{Аналитический регулятор траекторного слежения}

%В качестве нижнего модуля управления рассматривается алгоритм, описанный в работе~\cite{kim_algorithms}.
%Основной идеей является движение вдоль заданной опорной траектории с обеспечением выполнение двух условий: сохранение малой ошибки относительно траектории и поддержание заданной скорости продвижения вдоль неё.
%\begin{gather}
%    \text{dist}(y, S) \leq e_{max},\\
%    |\ || \dot y || - V^{\ast}| \leq \tilde V_{max},
%\end{gather}
%где S — целевая траектория, $V^\ast$ — параметрическая скорость движения вдоль неё, а $e_{\max}$ и $\tilde V_{\max}$ задают допустимые отклонения.
%На основе скорости $V^{\ast}$ формируется опорная точка $s_{ref}$.
%В диссертации~\cite{kim_algorithms} $s_{ref}$ задается выражением $s_{ref} = V^\ast t$, но это работает для постоянной скорости.
%В текущей постановке задачи разумнее всего вычислять $s_{ref}$ численным приближением:
%\begin{gather}
%    s_{ref} = \limit\int_0^t V^\ast (\tau) d\tau
%\end{gather}
%Это обусловлено изменчивостью параметра $V^{\ast}$ и при использовании подхода из диссертации~\cite{kim_algorithms} опорная точка имела бы скачкообразный характер, что для задачи слежения критически недопустимо.
%Далее на основе текущей опорной точки и положения объекта формируются ошибки слежения $e_1, e_2$ и задается управляющее воздействие.
%В результате такой последовательности действий объект корректирует скорость движения.

Пусть заданная траектория описывается параметрической кривой
\begin{gather}
    p_r(s) = [x_r(s), y_r(s), z_r(s)]^T,
\end{gather}
где $s$~--- параметр продвижения вдоль траектории. Для текущего положения
квадрокоптера $p(t)$ определяется ближайшая точка траектории
$p_r(s_{\mathrm{ref}})$, после чего ошибка слежения записывается как
\begin{gather}
    e_p = p(t) - p_r(s_{\mathrm{ref}}).
\end{gather}
Далее эта ошибка проецируется на локальный базис, связанный с траекторией.
В результате формируются продольная и поперечная составляющие ошибки:
\begin{gather}
    e_1 = e_p^T t_r, \qquad
    e_2 = \|e_p - e_1 t_r\|,
\end{gather}
где $t_r$~--- единичный касательный вектор к траектории в точке
$s_{\mathrm{ref}}$.

Задача регулятора состоит в формировании такого управляющего воздействия,
при котором ошибки $e_1$ и $e_2$ остаются ограниченными, а объект продолжает
движение вдоль траектории с заданной параметрической скоростью $V^*$.
В общем виде опорная скорость может быть записана как
\begin{gather}
    \dot s_{\mathrm{ref}} = V^*.
\end{gather}
При постоянном значении $V^*$ это даёт $s_{\mathrm{ref}} = V^*t$, однако
в данной работе скорость выбирается адаптивно, поэтому используется
интегральная форма:
\begin{gather}
    s_{\mathrm{ref}}(t)=\int_0^t V^*(\tau)d\tau.
\end{gather}

Управляющее воздействие нижнего уровня формируется по ошибкам слежения
и их производным:
\begin{gather}
    u = u(e_1, e_2, \dot e_1, \dot e_2, V^*).
\end{gather}
В рамках данной работы интеллектуальный модуль не изменяет структуру
аналитического регулятора, а влияет только на значение $V^*$. Поэтому
устойчивость движения обеспечивается базовым контуром слежения, а обучаемая
модель выполняет функцию верхнего уровня~--- адаптивно выбирает допустимую
скорость продвижения вдоль траектории.


\section{Формирование датасета}\label{sec:form_dataset}

Перед обучением моделей был сформирован датасет локальных состояний системы и максимально возможной скорости на этом участке.
Каждый пример создается в симуляции.
Сначала выбирается тип траектории и точка на ней, затем формируется возмущенное состояние квадрокоптера, после чего вычисляется набор признаков и через процедуру перебора оценивается максимально допустимое значение $V^{\ast}$, при котором базовый контур еще сохраняет качественное слежение.

В текущей реализации использован следующий вектор признаков:
\begin{gather}
    f =
    \begin{bmatrix}
        e_1 & e_2 & \dot{e}_2 & \lVert v \rVert & e_{\mathrm{head}} & \kappa_c & \kappa_{\max}
    \end{bmatrix}^{\top}.
\end{gather}
Здесь $e_1$ и $e_2$ описывают продольное и поперечное отклонения, $\dot{e}_2$ отражает динамику поперечной ошибки, $\lVert v \rVert$ соответствует текущей скорости движения, $e_{\mathrm{head}}$ характеризует угловое рассогласование с касательной, а $\kappa_c$ и $\kappa_{\max}$ описывают локальную кривизну траектории.

Такой выбор признаков строится по принципу содержательной достаточности.
Если объект уже имеет заметную поперечную ошибку, скорость лучше уменьшить.
Если впереди участок с большой кривизной, запас по скорости тоже должен снижаться.
Если же маршрут локально близок к прямой, а текущее движение стабильно, можно разрешить более высокий темп.


\section{Методы машинного обучения для задачи регрессии скорости}

Этот раздел посвящен описанию и введению базовых понятий и алгоритмов, которые будут использоваться далее в описании экспериментов и результатов.

\subsection{Обучение с учителем и регрессия}

В данной задаче имеется набор размеченных примеров: вектор признаков~$f_i$ и соответствующее целевое значение~$V_{\mathrm{opt},i}$, полученное от оракула.
Цель - обучить модель~$M_\theta$, которая по новому вектору признаков выдавала бы корректный прогноз скорости.
В данной работе задача выбора скорости рассматривается как задача регрессии.
По вектору признаков модель должна предсказать допустимое значение V_{opt}.
Для обучения используется среднеквадратичная ошибка.

\begin{gather}
    \mathcal{L}_{\mathrm{MSE}} = \frac{1}{N}\sum_{i=1}^N \left(\hat{V}_i - V_{\mathrm{opt},i}\right)^2.
\end{gather}

\subsection{Полносвязная нейронная сеть}

Полносвязная нейронная сеть (MLP) последовательно преобразует входной вектор через несколько слоёв. Каждый нейрон~$j$ слоя~$l$ выполняет операцию
\begin{gather}
    a_j^{(l)} = \sigma\!\left(\sum_k w_{jk}^{(l)}\, a_k^{(l-1)} + b_j^{(l)}\right),
\end{gather}
где $\sigma$ - нелинейная функция активации, $w_{jk}^{(l)}$ - обучаемые веса, $b_j^{(l)}$ - смещение. В данной работе используется ReLU:
\begin{gather}
    \mathrm{ReLU}(z) = \max(0,\, z).
\end{gather}
Чередование линейных преобразований и нелинейностей позволяет сети аппроксимировать сложные зависимости.

\subsection{Градиентный спуск и оптимизатор}

Параметры~$\theta$ подбираются итеративной минимизацией функции потерь~$\mathcal{L}(\theta)$.
В данной работе используется оптимизатор Adam~\cite{kingma2015adam}, который адаптивно масштабирует шаг для каждого параметра:
\begin{gather}
    \theta_{t+1} = \theta_t - \frac{\eta}{\sqrt{\hat{v}_t} + \varepsilon}\,\hat{m}_t,
\end{gather}
где $\hat{m}_t$ и $\hat{v}_t$ - адаптивные оценки первого и второго моментов градиента, $\eta$ - скорость обучения, $\varepsilon$ - предотвращение деления на ноль.\\
Также для контроля переобучения данные разделялись на обучающую и валидационную выборки.
Обучение прекращалось при отсутствии улучшения ошибки на валидации.



%\section{Модель MLP и почему она нужна}
%
%Базовой (baseline) моделью в проекте является полносвязная нейронная сеть.
%Так как она достаточна понятна в реализации - на вход подается небольшой набор числовых признаков, затем несколько скрытых слоев последовательно преобразуют их.
%И на выходе получается одно число - прогноз допустимой скорости.


\section{Модели обучения с подкреплением}

Помимо MLP в работе исследованы три архитектуры, основанные на идеях обучения с подкреплением: SAC, TD3 и PPO.
В классической постановке эти методы учатся выбирать действия в среде, максимизируя накопленную награду.
В данном проекте их идея адаптирована под задачу предсказания скорости: модели получают тот же вектор признаков и должны выдать одно скалярное действие - значение $V^{\ast}$.
В общем виде политика выбора описывается как:
\begin{gather}
    \pi(a|s) = \arg\max\mathbb E \left[ \limit\sum_{t=0}^{\infty}\gamma^t r_t \right],
\end{gather}
где $a$ - действие, $s$ - состояние.
В классическом понимании алгоритмов обучения с подкреплением под действием понимается управляющее воздействие на объект.
Но в данной работе в роли действия выступает параметрическая скорость движения $V^\ast$.
А состояние формируется из вектора признаков.
И тогда политика формируется как:
\begin{gather}
    V^{\ast} = \pi(s)
\end{gather}

\textbf{SAC.} Метод Soft Actor-Critic использует стохастическую политику и стремится не только максимизировать полезность действия, но и сохранять достаточную энтропию распределения~\cite{haarnoja2018sac}.
\begin{gather}
    J(\pi) = \mathbb{E}\left[ \limit\sum_t (r_t + \alpha H(\pi(\dot | s_t))) \right], \\
    H(\pi) = -\mathbb{E}[\log{\pi(a|s)}] \text{ - энтропия политики}
\end{gather}
%Модель не фиксируется на одной стратегии, а продолжает аккуратно подбирать возможное лучшее действие.
%Для задачи выбора скорости это полезно тем, что SAC обычно учится более плавному и устойчивому поведению на разных режимах.

\textbf{TD3.} Алгоритм Twin Delayed DDPG строится вокруг детерминированной политики и пары критиков, которые уменьшают переоценку качества действий~\cite{fujimoto2018td3}.
\begin{gather}
    y = r + \gamma \min(Q_1\textasciiacute, Q_2\textasciiacute)
\end{gather}
%На интуитивном уровне использование двух критиков уменьшает риск переоценки качества отдельных действий.
%Поэтому TD3 часто стремится к более агрессивным решениям, если они дают выигрыш по целевой функции.
%В задачах управления это может быть плюсом по скорости, но повышает риск нестабильного поведения на сложных сценариях.

\textbf{PPO.} Метод Proximal Policy Optimization использует ограниченное обновление политики, чтобы обучение не делало слишком резких шагов~\cite{schulman2017ppo}.
\begin{gather}
    L^{\text{CLIP}} = \mathbb E [\min(r_t(\theta)A_t, clip(r_t(\theta), 1 - \varepsilon, 1 + \varepsilon)A_t)],\\
    clip =
    \begin{cases}
        a, x < a\\
        x, a \leq x \leq b\\
        b, b < x
    \end{cases}.
\end{gather}
За счет этого PPO предотвращает резкие изменения политики между итерациями.
%Это дает более консервативное изменение стратегии и является компромиссом между скоростью обучения и устойчивостью результата.
%Но для достижения максимального результата следует использовать тонкую настройку функцию награды.


\section{Метрики качества моделей}

Для дополнительной оценки использовался коэффициент детерминации $R^2$, характеризующий долю объяснённой дисперсии целевой переменной.
\begin{gather}
    R^2 = 1 - \frac{\limit\sum_i \left(\hat{V}_i - V_{\text{opt},i}\right)^2}{\limit\sum_i \left( \bar{V} - V_{\text{opt},i} \right)^2},
\end{gather}


\section{Интеграция нейросетевого модуля в архитектуру управления}

Имея верхний и нижний уровень описания модели квадрокоптера, необходимо было их соединить для получения единой рабочей системы.
Общий принцип работы итерационно описывается примерно так.
На каждом шаге моделирования вычисляется вектор признаков на основе текущего состояния системы.
Далее модель предсказывает новое значение $V^\ast$ и уже это значение пропускается через набор условных ограничений для корректной работы программы.
И наконец подается в базовый регулятор, который применяет воздействие на систему.
Для корректной работы пришлось внести в аналитический регулятор ряд ограничений, которые решают проблемы расхождения.
\begin{enumerate}
    \item Был введен дополнительный контрольный регулятор, который обеспечивает плавность изменения скорости, для предотвращения скачков и соблюдения синхронизации.
    \item Предсказанное значение должно находиться в допустимом диапазоне скоростей для конкретной модели квадрокоптера, дабы исключить физически невозможные значения.
    \item На старте моделирования некоторое время алгоритм адаптации находится в деактивированном состоянии, чтобы система успела стабилизироваться около траектории.
\end{enumerate}

Именно этот этап превратил набор отдельных уровней работы в полноценное рабочее решение.
%
%\section{Технические проблемы интеграции и их решения}
%\label{sec:integration_issues}
%
%Подключение нейросетевого модуля к аналитическому контуру выявило ряд нетривиальных инженерных проблем. Каждая из них неочевидна при разработке модели или регулятора по отдельности, но хорошо заметна в замкнутой симуляции.
%
%\textbf{Разрыв опорного сигнала при изменении скорости.} Исходная реализация использовала $s_{\mathrm{ref}} = V^{\ast} \cdot t$ для первой компоненты ошибки слежения. При изменении $V^{\ast}$ эта формула порождала мгновенный скачок, который наблюдатель высокого усиления воспринимал как резкий импульс. Решение - интегральное накопление $s_{\mathrm{ref}} = \int_0^t V^{\ast}(\tau)\,d\tau$, которое гарантирует непрерывность при любом законе изменения скорости.
%
%\textbf{Накопление $\eta$-интегратора (windup).} При высоком $V^{\ast}$ ошибки не убывали достаточно быстро, и~$\eta$-интегратор накапливал большое значение, приводя к взрывному росту управляющего сигнала. Введён монитор безопасности: при $|e_2| > 0{.}5 \cdot e_{2,\mathrm{lim}}$ или при $s_{\mathrm{arc}} - s_{\mathrm{ref}} < -0{.}3$~м нейросетевому модулю запрещается увеличивать $V^{\ast}$.
%
%\textbf{Двойной вызов \texttt{step()} при инициализации.} Цикл симуляции вызывал шаг контроллера дважды при $t=0$ (pre-loop), из-за чего состояние наблюдателя загрязнялось до начала основного цикла. После первого (pre-loop) вызова состояние контроллера сбрасывается, и основной цикл стартует из чистого состояния.
%
%\textbf{Нестабильные предсказания при малом датасете.} При ограниченном объёме обучающей выборки модели порой предсказывали $V^{\ast}$ вблизи верхней границы диапазона независимо от текущей геометрии траектории. Комбинация верхнего ограничения скорости (\texttt{vstar\_cap}) и ограничителя темпа изменения (\texttt{vstar\_rate}) позволила сгладить такое поведение и обеспечить безопасный режим работы.

%\section{Модуль автоматизированной оценки}
%
%Чтобы сравнение моделей было воспроизводимым, в проекте реализован отдельный модуль автоматизированной оценки. Он запускает одинаковые сценарии для baseline-режима и всех обученных моделей, собирает одинаковый набор метрик и строит унифицированные графики. Это важно не только для удобства. Если каждая модель тестируется в слегка разных условиях, само сравнение теряет смысл.
%
%В проекте модуль оценки строит временны\'е зависимости ошибок и скорости, формирует сводные диаграммы и автоматически готовит таблицу результатов для включения в отчет. За счет этого можно сравнивать модели не по отдельным удачным примерам, а по единому набору бенчмарков.


\section{Выводы по главе}

Во второй главе описана теоретическая часть построения интеллектуальной системы решения.
Она включает описание работы аналитического контура комплекса, геометрический раздел слежения за кривой, процедуру формирования данных, раздел машинного обучения с пояснением архитектур и саму реализацию связи нейросетевого модуля с аналитическим.
Также можно выделить решение ряда технических проблем, которые возникли в ходе интеграции.
В качестве задатка для дальнейшего исследования следует отметить, что целесообразно расширить учитываемые показатели числом резких изменений $V^\ast$, энергетическими затратами на маневры и устойчивостью к геометрии.
Чего не было учтено в текущей реализации, так как эти метрики особенно важны при переносе системы из симуляции в физическую среду.